\small{
\noindent 
Dynamic epidemic models have proven valuable for public health decision makers as they provide useful insights into the understanding and prevention of infectious diseases. However, inference for these types of models can be difficult because the disease spread is typically only partially observed e.g. in form of reported incidences in given time periods. This chapter discusses how to perform likelihood-based inference for partially observed Markov epidemic models when it is relatively easy to generate samples from the Markov transmission model while the likelihood function is intractable. The first part of the chapter reviews the theoretical background of inference for partially observed Markov processes (POMP) via iterated filtering. In the second part of the chapter the performance of the method and associated practical difficulties are illustrated on two examples. In the first example a simulated outbreak data set consisting of the number of newly reported cases aggregated by week is fitted to a POMP where the underlying disease transmission model is assumed to be a simple Markovian SIR model. The second example illustrates possible model extensions such as seasonal forcing and over-dispersion in both, the transmission and observation model, which can be used, e.g., when analysing routinely collected rotavirus surveillance data. Both examples are implemented using the R-package \texttt{pomp}  \citep{kingNgu} and the code is made available online. 
}
